%% =====================================================================
%%  T-02-02  Indifference to Being Liked
%%  -------------------------------------------------------------------
%%  One idea per file. Core is mandatory. Slots in canonical order.
%%  Max 700 words. No layout commands. No filler. New source material
%%  for this entry goes in sources/<BOOK>/T-02-02.tex -- never by
%%  rewriting this file (docs/RULES.md R0.1).
%% =====================================================================
%% @kind: topic
%% @id: T-02-02
%% @title: Indifference to Being Liked
%% @subtitle: The willingness to be disliked is the precondition
%% @chapter: c02
%% @order: 20
%% @status: stable
%% @sources: B1, B3
%% @updated: 2026-09-19
%% @allow-rewrite: slot migration to the seven-question format (docs/RULES.md section 2)

\topic{T-02-02}{Indifference to Being Liked}{The willingness to be disliked is the precondition}


\begin{core}
Most people lose exchanges because being liked is worth more to them than the thing being negotiated. The practitioner has decided otherwise, at least for the duration of the exchange, and that single difference accounts for a large part of the gap described in T-01-01.
\end{core}

\begin{mechanism}
Being liked is a real asset and treating it as one is correct; the error is treating it as an asset that must never be spent. A person unwilling to spend it telegraphs that in advance, and the counterpart prices it in. Two consequences follow. First, the unwilling person cannot refuse, cannot delay, and cannot let an accusation stand unrefuted --- which removes most of the defences in Part VI. Second, they over-explain, and over-explaining is itself a tell. The corollary is that retaliation is almost always a bad trade: winning an argument buys a moment and creates an enemy with a long memory.
\end{mechanism}

\begin{application}
  \item Before a difficult exchange, decide in advance what you are willing to be disliked for.
  \item Answer the question asked and stop. Do not fill the silence that follows.
  \item Let small misreadings of you stand uncorrected; spend the correction only where it changes an outcome.
  \item Refuse retaliation as a category: it converts a settled matter into an ongoing one.
\end{application}

\begin{feedback}
  \item You explain yourself more than the situation requires.
  \item You have conceded something in order to end an uncomfortable silence.
  \item You are managing how you are perceived by someone whose opinion has no bearing on anything you want.
  \item You let a false statement about you stand because correcting it would be awkward.
\end{feedback}

\begin{failure}
Carried too far this produces the person nobody will help, and both sources are explicit that being needed and being trusted are assets worth holding. The usable version is selective: indifference to being liked by the counterpart in a contested exchange, not by everyone. The source that argues for meekness argues for it on effectiveness grounds --- retaliation escalates and costs more than it recovers --- not on grounds of contempt for the other person.
\end{failure}

\begin{countermeasures}
  \item Notice when warmth is being purchased rather than given. A counterpart who works hard to be liked is usually about to ask for something.
  \item Do not reward over-explanation with more pressure; it trains the other person to keep conceding.
  \item If you are the one who cannot bear to be disliked, put the decision in writing or hand it to someone else --- a rule, a policy, a third party --- so that the refusal is not personally yours to make.
\end{countermeasures}

\sources{B1, B3}
\seesources
\seealso{T-02-01, T-06-01, T-16-03, T-26-02}
